\ifdefined\gkver\else\def\gkver{student}\fi
\documentclass[\gkver]{gaokaozhenti}
\begin{document}

\chapter{2015年北京理综}
%% paperType: 理综物理部分

\begin{enumerate}
\item
%% number: 13
%% paperName: 2015年北京理综第 13 题 6 分
%% typeId: 1
%% type: 单选题
%% chapter: 气体性质
%% point: 物体的内能
%% method: 无
%% score: 6
%% degree: 650
%% duplicateId: 0
%% body:
下列说法正确的是\xzanswer{C}
\fourchoices[answer=C]
{物体放出热量，其内能一定减小}
{物体对外做功，其内能一定减小}
{物体吸收热量，同时对外做功，其内能可能增加}
{物体放出热量，同时对外做功，其内能可能不变}
%% answer: C
%% memo:
\memoanswer{
	本题原卷及材料中未提供详解，待补充。
}

\item
%% number: 14
%% paperName: 2015年北京理综第 14 题 6 分
%% typeId: 1
%% type: 单选题
%% chapter: 原子和原子核
%% point: 天然放射性
%% method: 无
%% score: 6
%% degree: 650
%% duplicateId: 0
%% body:
下列核反应方程中，属于 $\alpha$ 衰变的是\xzanswer{B}
\fourchoices[answer=B]
{\ce{^{14}_{7}N + ^{4}_{2}He -> ^{17}_{8}O + ^{1}_{1}H}}
{\ce{^{238}_{92}U -> ^{234}_{90}Th + ^{4}_{2}He}}
{\ce{^{2}_{1}H + ^{3}_{1}H -> ^{4}_{2}He + ^{1}_{0}n}}
{\ce{^{234}_{90}Th -> ^{234}_{91}Pa + ^{0}_{-1}e}}
%% answer: B
%% memo:
\memoanswer{
	本题原卷及材料中未提供详解，待补充。
}

\item
%% number: 15
%% paperName: 2015年北京理综第 15 题 6 分
%% typeId: 1
%% type: 单选题
%% chapter: 机械波
%% point: 波的图象
%% method: 无
%% score: 6
%% degree: 650
%% duplicateId: 0
%% body:
周期为 $2.0\Us$ 的简谐横波沿 $x$ 轴传播，该波在某时刻的图像如图所示，此时质点 P 沿 $y$ 轴负方向运动。则该波\xzanswer{B}

\onepicture[h=3.2cm]{figs/15.png}

\fourchoices[answer=B]
{沿 $x$ 轴正方向传播，波速 $v=20\Ums$}
{沿 $x$ 轴正方向传播，波速 $v=10\Ums$}
{沿 $x$ 轴负方向传播，波速 $v=20\Ums$}
{沿 $x$ 轴负方向传播，波速 $v=10\Ums$}
%% answer: B
%% memo:
\memoanswer{
	本题原卷及材料中未提供详解，待补充。
}

\item
%% number: 16
%% paperName: 2015年北京理综第 16 题 6 分
%% typeId: 1
%% type: 单选题
%% chapter: 曲线运动
%% point: 天体运动
%% method: 无
%% score: 6
%% degree: 650
%% duplicateId: 0
%% body:
假设地球和火星都绕太阳做匀速圆周运动，已知地球到太阳的距离小于火星到太阳的距离，那么\xzanswer{D}
\fourchoices[answer=D]
{地球公转周期大于火星的公转周期}
{地球公转的线速度小于火星公转的线速度}
{地球公转的加速度小于火星公转的加速度}
{地球公转的角速度大于火星公转的角速度}
%% answer: D
%% memo:
\memoanswer{
	本题原卷及材料中未提供详解，待补充。
}

\item
%% number: 17
%% paperName: 2015年北京理综第 17 题 6 分
%% typeId: 1
%% type: 单选题
%% chapter: 磁场
%% point: 带电粒子在磁场中的运动
%% method: 无
%% score: 6
%% degree: 650
%% duplicateId: 0
%% body:
实验观察到，静止在匀强磁场中 A 点的原子核发生 $\beta$ 衰变，衰变产生的新核与电子恰在纸面内做匀速圆周运动，运动方向和轨迹示意如图。则\xzanswer{D}

\onepicture[h=3.0cm]{figs/17.png}

\fourchoices[answer=D]
{轨迹 1 是电子的，磁场方向垂直纸面向外}
{轨迹 2 是电子的，磁场方向垂直纸面向外}
{轨迹 1 是新核的，磁场方向垂直纸面向里}
{轨迹 2 是新核的，磁场方向垂直纸面向里}
%% answer: D
%% memo:
\memoanswer{
	本题原卷及材料中未提供详解，待补充。
}

\item
%% number: 18
%% paperName: 2015年北京理综第 18 题 6 分
%% typeId: 1
%% type: 单选题
%% chapter: 运动和力
%% point: 变加速直线运动
%% method: 无
%% score: 6
%% degree: 450
%% duplicateId: 0
%% body:
“蹦极”运动中，长弹性绳的一端固定，另一端绑在人身上，人从几十米高处跳下。将蹦极过程简化为人沿竖直方向的运动。从绳恰好伸直，到人第一次下降至最低点的过程中，下列分析正确的是\xzanswer{A}
\fourchoices[answer=A]
{绳对人的冲量始终向上，人的动量先增大后减小}
{绳对人的拉力始终做负功，人的动能一直减小}
{绳恰好伸直时，绳的弹性势能为零，人的动能最大}
{人在最低点时，绳对人的拉力等于人所受的重力}
%% answer: A
%% memo:
\memoanswer{
	本题原卷及材料中未提供详解，待补充。
}

\item
%% number: 19
%% paperName: 2015年北京理综第 19 题 6 分
%% typeId: 1
%% type: 单选题
%% chapter: 电路
%% point: 分压与分流
%% method: 无
%% score: 6
%% degree: 450
%% duplicateId: 0
%% body:
如图所示，其中电流表 A 的量程为 $0.6\UA$，表盘均匀划分为 30 个小格，每一小格表示 $0.02\UA$；$R_{1}$ 的阻值等于电流表内阻的 $\frac{1}{2}$；$R_{2}$ 的阻值等于电流表内阻的 2 倍。若用电流表 A 的表盘刻度表示流过接线柱 1 的电流值，则下列分析正确的是\xzanswer{C}

\onepicture[h=3.2cm]{figs/19.png}

\fourchoices[answer=C]
{将接线柱 1、2 接入电路时，每一小格表示 $0.04\UA$}
{将接线柱 1、2 接入电路时，每一小格表示 $0.02\UA$}
{将接线柱 1、3 接入电路时，每一小格表示 $0.06\UA$}
{将接线柱 1、3 接入电路时，每一小格表示 $0.01\UA$}
%% answer: C
%% memo:
\memoanswer{
	本题原卷及材料中未提供详解，待补充。
}

\item
%% number: 20
%% paperName: 2015年北京理综第 20 题 6 分
%% typeId: 1
%% type: 单选题
%% chapter: 交流电
%% point: LC振荡回路
%% method: 无
%% score: 6
%% degree: 450
%% duplicateId: 0
%% body:
利用所学物理知识，可以初步了解常用的公交一卡通（IC 卡）的工作原理及相关问题。IC 卡内部有一个由电感线圈 L 和电容 C 构成的 LC 振荡电路。公交车上的读卡机（刷卡时“嘀”的响一声的机器）向外发射某一特定频率的电磁波。刷卡时，IC 卡内的线圈 $L$ 中产生感应电流，给电容 C 充电，达到一定的电压后，驱动卡内芯片进行数据处理和传输。下列说法正确的是\xzanswer{B}
\fourchoices[answer=B]
{IC 卡工作所需要的能量来源于卡内的电池}
{仅当读卡机发射该特定频率的电磁波时，IC 卡才能有效工作}
{若读卡机发射的电磁波偏离该特定频率，则线圈 $L$ 中不会产生感应电流}
{IC 卡只能接收读卡机发射的电磁波，而不能向读卡机传输自身的数据信息}
%% answer: B
%% memo:
\memoanswer{
	本题原卷及材料中未提供详解，待补充。
}

\item
%% number: 21
%% paperName: 2015年北京理综第 21 题 15 分
%% typeId: 4
%% type: 实验题
%% chapter: 机械振动
%% point: 单摆测重力加速度
%% method: 无
%% score: 15
%% degree: 450
%% duplicateId: 0
%% body:
（1）“测定玻璃的折射率”的实验中，在白纸上放好玻璃砖，aa$'$ 和 bb$'$ 分别是玻璃砖与空气的两个界面，如图所示。在玻璃砖的一侧插上两枚大头针 P$_{1}$ 和 P$_{2}$，用“$+$”表示大头针的位置，然后在另一侧透过玻璃砖观察，并依次插上大头针 P$_{3}$ 和 P$_{4}$。在插 P$_{3}$ 和 P$_{4}$ 时，应使\tkanswer[1.6cm]{C}（选填选项前的字母）。

（A）P$_{3}$ 只挡住 P$_{1}$ 的像

（B）P$_{4}$ 只挡住 P$_{2}$ 的像

（C）P$_{3}$ 同时挡住 P$_{1}$、P$_{2}$ 的像

\onepicture[h=2.8cm, align=right]{figs/21.png}

（2）用单摆测定重力加速度的实验装置如图 1 所示。

\begin{steps}
	\item
	组装单摆时，应在下列器材中选用\tkanswer[1.6cm]{A、D}（选填选项前的字母）。

	（A）长度为 $1\Um$ 左右的细线

	（B）长度为 $30\Ucm$ 左右的细线

	（C）直径为 $1.8\Ucm$ 的塑料球

	（D）直径为 $1.8\Ucm$ 的铁球

	\item
	测出悬点 O 到小球球心的距离（摆长）$L$ 及单摆完成 $n$ 次全振动所用的时间 $t$，则重力加速度 $g=$\tkanswer[2.4cm]{$\frac{4\pi^{2}n^{2}L}{t^{2}}$}（用 $L$、$n$、$t$ 表示）。

	\item
	下表是某同学记录的 3 组实验数据，并做了部分计算处理。

	\begin{center}
	\begin{tabular}{|c|c|c|c|}
	\hline
	组次 & 1 & 2 & 3 \\
	\hline
	$L/\Ucm$ & 80.00 & 90.00 & 100.00 \\
	\hline
	$t/\Us$ & 90.0 & 95.5 & 100.5 \\
	\hline
	$T/\Us$ & 1.80 & 1.91 & \\
	\hline
	$g$（$\Umsq$） & 9.74 & 9.73 & \\
	\hline
	\end{tabular}
	\end{center}

	请计算出第 3 组实验中的 $T=$\tkanswer[1.2cm]{2.01}$\Us$，$g=$\tkanswer[1.2cm]{9.76}$\Umsq$。

	\item
	用多组实验数据做出 $T^{2}-L$ 图像，也可以求出重力加速度 $g$。已知三位同学做出的 $T^{2}-L$ 图线的示意图如图 2 中的 a、b、c 所示，其中 a 和 b 平行，b 和 c 都过原点，图线 b 对应的 $g$ 值最接近当地重力加速度的值。则相对于图线 b，下列分析正确的是\tkanswer[1.6cm]{B}（选填选项前的字母）。

	（A）出现图线 a 的原因可能是误将悬点到小球下端的距离记为摆长 $L$

	（B）出现图线 c 的原因可能是误将 49 次全振动记为 50 次

	（C）图线 c 对应的 $g$ 值小于图线 b 对应的 $g$ 值

	\item
	某同学在家里测重力加速度。他找到细线和铁锁，制成一个单摆，如图 3 所示。由于家里只有一根量程为 $30\Ucm$ 的刻度尺，于是他在细线上的 A 点做了一个标记，使得悬点 O 到 A 点间的细线长度小于刻度尺量程。保持该标记以下的细线长度不变，通过改变 O、A 间细线长度以改变摆长。实验中，当 O、A 间细线的长度分别为 $l_{1}$、$l_{2}$ 时，测得相应单摆的周期为 $T_{1}$、$T_{2}$。由此可得重力加速度 $g=$\tkanswer[2.6cm]{$\frac{4\pi^{2}(l_{1}-l_{2})}{T_{1}^{2}-T_{2}^{2}}$}（用 $l_{1}$、$l_{2}$、$T_{1}$、$T_{2}$ 表示）。
\end{steps}

\threepicture[showlabel=false, align=right, hA=4.2cm, hB=3.0cm, hC=3.4cm]
{figs/21a.png}
{figs/21b.png}
{figs/21c.png}

%% answer:
\jdanswer{
\begin{enumerate}
	\item
	C
	\item
	①A、D；②$\frac{4\pi^{2}n^{2}L}{t^{2}}$；③2.01，9.76；④B；⑤$\frac{4\pi^{2}(l_{1}-l_{2})}{T_{1}^{2}-T_{2}^{2}}$
\end{enumerate}
}
%% memo:
\memoanswer{
	本题原卷及材料中未提供详解，待补充。
}

\item
%% number: 22
%% paperName: 2015年北京理综第 22 题 16 分
%% typeId: 6
%% type: 计算题
%% chapter: 电磁感应
%% point: 电磁感应电路分析
%% method: 无
%% score: 16
%% degree: 650
%% duplicateId: 0
%% body:
如图所示，足够长的平行光滑金属导轨水平放置，宽度 $L=0.4\Um$，一端连接 $R=1\UO$ 的电阻。导轨所在空间存在竖直向下的匀强磁场，磁感应强度 $B=1\UT$。导体棒 MN 放在导轨上，其长度恰好等于导轨间距，与导轨接触良好。导轨和导体棒的电阻均可忽略不计。在平行于导轨的拉力 $F$ 作用下，导体棒沿导轨向右匀速运动，速度 $v=5\Ums$。求：

\begin{enumerate}
	\item
	感应电动势 $E$ 和感应电流 $I$；
	\item
	在 $0.1\Us$ 时间内，拉力的冲量 $I_{F}$ 的大小；
	\item
	若将 MN 换为电阻 $r=1\UO$ 的导体棒，其它条件不变，求导体棒两端的电压 $U$。
\end{enumerate}

\onepicture[h=3.2cm, align=right]{figs/22.png}

%% answer:
\jdanswer{
\begin{enumerate}
	\item
	$E=2\UV$，$I=2\UA$
	\item
	$I_{F}=0.08\UNs$
	\item
	$U=1\UV$
\end{enumerate}
}
%% memo:
\memoanswer{
	（1）由法拉第电磁感应定律可得，感应电动势 $E=BLv=1\times0.4\times5\UV=2\UV$

	感应电流 $I=\frac{E}{R}=\frac{2}{1}\UA=2\UA$

	（2）拉力大小等于安培力大小 $F=BIL=1\times2\times0.4\UN=0.8\UN$

	冲量大小 $I_{F}=F\Delta t=0.8\times0.1\UNs=0.08\UNs$

	（3）由闭合电路欧姆定律可得，电路中电流 $I'=\frac{E}{R+r}=\frac{2}{2}\UA=1\UA$

	由欧姆定律可得，导体棒两端电压 $U=I'R=1\UV$
}

\item
%% number: 23
%% paperName: 2015年北京理综第 23 题 18 分
%% typeId: 6
%% type: 计算题
%% chapter: 功和能
%% point: 动能定理
%% method: 无
%% score: 18
%% degree: 450
%% duplicateId: 0
%% body:
如图所示，弹簧的一端固定，另一端连接一个物块，弹簧质量不计。物块（可视为质点）的质量为 $m$，在水平桌面上沿 $x$ 轴运动，与桌面间的动摩擦因数为 $\mu$。以弹簧原长时物块的位置为坐标原点 $O$，当弹簧的伸长量为 $x$ 时，物块所受弹簧弹力大小为 $F=kx$，$k$ 为常量。

\begin{enumerate}
	\item
	请画出 $F$ 随 $x$ 变化的示意图；并根据 $F-x$ 图像求物块沿 $x$ 轴从 $O$ 点运动到位置 $x$ 的过程中弹力所做的功。

	\item
	物块由 $x_{1}$ 向右运动到 $x_{3}$，然后由 $x_{3}$ 返回到 $x_{2}$，在这个过程中，

	（A）求弹力所做的功，并据此求弹性势能的变化量；

	（B）求滑动摩擦力所做的功；并与弹力做功比较，说明为什么不存在与摩擦力对应的“摩擦力势能”的概念。
\end{enumerate}

\onepicture[h=2.4cm, align=right]{figs/23.png}

%% answer:
\jdanswer{
\begin{enumerate}
	\item
	$W=-\frac{1}{2}kx^{2}$
	\item
	（A）$\Delta E_{p}=\frac{1}{2}kx_{2}^{2}-\frac{1}{2}kx_{1}^{2}$；（B）$W_{f}=-\mu mg(2x_{3}-x_{1}-x_{2})$，不存在与摩擦力对应的“摩擦力势能”
\end{enumerate}
}
%% memo:
\memoanswer{
	（1）$F-x$ 图像如图。

	物块沿 $x$ 轴从 $O$ 点运动到位置 $x$ 的过程中，弹力做负功；$F-x$ 图线下的面积等于弹力做功大小，弹力做功

	$W=-\frac{1}{2}kx\cdot x=-\frac{1}{2}kx^{2}$

	\onepicture[h=3.2cm, align=right]{figs/23b.jpg}

	（2）（A）物块由 $x_{1}$ 向右运动到 $x_{3}$ 的过程中，弹力做功

	$W_{T1}=-\frac{1}{2}(kx_{1}+kx_{3})(x_{3}-x_{1})=\frac{1}{2}kx_{1}^{2}-\frac{1}{2}kx_{3}^{2}$

	物块由 $x_{3}$ 向左运动到 $x_{2}$ 的过程中，弹力做功

	$W_{T2}=-\frac{1}{2}(kx_{2}+kx_{3})(x_{3}-x_{2})=\frac{1}{2}kx_{3}^{2}-\frac{1}{2}kx_{2}^{2}$

	整个过程中，弹力做功

	$W_{T}=W_{T1}+W_{T2}=\frac{1}{2}kx_{1}^{2}-\frac{1}{2}kx_{2}^{2}$

	弹性势能的变化量

	$\Delta E_{p}=-W_{T}=\frac{1}{2}kx_{2}^{2}-\frac{1}{2}kx_{1}^{2}$

	（B）整个过程中，摩擦力做功

	$W_{f}=-\mu mg(2x_{3}-x_{1}-x_{2})$

	与弹力做功比较，弹力做功与 $x_{3}$ 无关，即与实际路径无关，只与始末位置有关，所以，我们可以定义一个由物体之间的相互作用力（弹力）和相对位置决定的能量——弹性势能，而摩擦力做功与 $x$ 有关，即与实际路径有关，所以，不可以定义与摩擦力对应的“摩擦力势能”。
}

\item
%% number: 24
%% paperName: 2015年北京理综第 24 题 20 分
%% typeId: 6
%% type: 计算题
%% chapter: 电场
%% point: 带电粒子的直线运动
%% method: 无
%% score: 20
%% degree: 250
%% duplicateId: 0
%% body:
真空中放置的平行金属板可以用作光电转换装置，如图所示。光照前两板都不带电。以光照射 A 板，则板中的电子可能吸收光的能量而逸出。假设所有逸出的电子都垂直于 A 板向 B 板运动，忽略电子之间的相互作用。保持光照条件不变。a 和 b 为接线柱。

已知单位时间内从 A 板逸出的电子数为 $N$，电子逸出时的最大动能为 $E_{km}$。元电荷为 $e$。

\begin{enumerate}
	\item
	求 A 板和 B 板之间的最大电势差 $U_{m}$，以及将 a、b 短接时回路中的电流 $I_{\text{短}}$。
	\item
	图示装置可看作直流电源，求其电动势 $E$ 和内阻 $r$。
	\item
	在 a 和 b 之间连接一个外电阻时，该电阻两端的电压为 $U$。外电阻上消耗的电功率设为 $P$；单位时间内到达 B 板的电子，在从 A 板运动到 B 板的过程中损失的动能之和设为 $\Delta E_{k}$。请推导证明：$P=\Delta E_{k}$。（注意：解题过程中需要用到、但题目没有给出的物理量，要在解题中做必要的说明）
\end{enumerate}

\onepicture[h=3.4cm, align=right]{figs/24.png}

%% answer:
\jdanswer{
\begin{enumerate}
	\item
	$U_{m}=\frac{E_{km}}{e}$，$I_{\text{短}}=Ne$
	\item
	$E=\frac{E_{km}}{e}$，$r=\frac{E_{km}}{Ne^{2}}$
	\item
	证明过程见详解，$P=\Delta E_{k}$
\end{enumerate}
}
%% memo:
\memoanswer{
	（1）由动能定理得：$E_{km}=eU$，可得 $U_{m}=\frac{E_{km}}{e}$

	短路时所有逸出电子都到达 B 板，故短路电流 $I_{\text{短}}=Ne$

	（2）电源的电动势等于断路时的路端电压，即上面求出的 $U_{m}$，所以 $E=U_{m}=\frac{E_{km}}{e}$

	电源内阻：$r=\frac{E}{I_{\text{短}}}=\frac{E_{km}}{Ne^{2}}$

	（3）外电阻两端的电压为 $U$，则电源两端的电压也是 $U$。

	由动能定理，一个电子经电源内部电场后损失的动能之和 $\Delta E_{\text{ke}}=eU$

	设单位时间内有 $N'$ 个电子到达 B 板，则损失的动能之和 $\Delta E_{k}=N'\Delta E_{\text{ke}}=N'eU$

	根据电流的定义，此时电源内部的电流 $I=N'e$，

	此时流过外电阻的电流也是 $I=N'e$，外电阻上消耗的电功率 $P=IU=N'eU$

	所以 $P=\Delta E_{k}$
}

\end{enumerate}

\end{document}
